\begin{solapartat}
\punts{0,25} $E_{\text{fotó}} = hf = h\dfrac{c}{\lambda}$, com més petita sigui $\lambda$,
major serà la freqüència i l'energia del fotó, per tant, l'ona de $\lambda = 550$ nm serà
la d'energia més gran.

Per al vermell:

\punts{0,1} $f_r = \dfrac{c}{\lambda_r} = 4,00\times 10^{14}$ Hz

\punts{0,2} $E_{\text{fotó},r} = hf_r = 2,65\times 10^{-19}$ J.

\punts{0,2} $E_{1\mathrm{mol},r} = N_A hf_r = 1,60\times 10^{5}$ J $= 160$ kJ.

Per al verd:

\punts{0,1} $f_g = \dfrac{c}{\lambda_g} = 5,45\times 10^{14}$ Hz

\punts{0,2} $E_{\text{fotó},g} = hf_g = 3,62\times 10^{-19}$ J.

\punts{0,2} $E_{1\mathrm{mol},g} = N_A hf_g = 2,18\times 10^{5}$ J $= 218$ kJ.
\end{solapartat}

\begin{solapartat}
\punts{0,15} La freqüència llindar correspon a la dels fotons que tenen una energia
$W_0$:
\[ f_0 = \frac{W_0}{h} = \frac{3,46\times 10^{-19}}{6,63\times 10^{-34}} = 5,22\times 10^{14}\ \mathrm{Hz} \]

\punts{0,1} $\lambda_0 = \dfrac{c}{f_0} = 5,75\times 10^{-7}$ m $= 575$ nm

% DUBTE: a l'original els exponents d'aquestes dues freqüències surten com "10^{15}"
% (4,00×10^15 Hz i 5,45×10^15 Hz), en comptes de 10^{14} com a l'apartat a); es transcriu
% tal com apareix a la pauta oficial.
\punts{0,1} Com que per al vermell $f_r = 4,00\times 10^{15}$ Hz $< f_0$, no té suficient
energia per arrencar electrons.

\punts{0,1} Com que per al verd $f_g = \dfrac{c}{\lambda_g} = 5,45\times 10^{15}$ Hz
$> f_0$, sí té suficient energia per arrencar electrons.

\punts{0,3} $E_{C,\text{màx}} = hf_g - W_0 = 1,56\times 10^{-20}$ J.

\punts{0,5}
\figura[0.55]{sol_fig1.png}

Hi ha dues regions:
\begin{itemize}
\item si l'energia del fotó és inferior a $W_0$ no s'arrenquen electrons, llavors $E_{C,\text{màx}} = 0$.
\item si l'energia del fotó és superior a $W_0$ s'arrenquen electrons i $E_{C,\text{màx}} = E_{\text{fotó}} - W_0$, és a dir, la dependència de $E_{C,\text{màx}}$ en $E_{\text{fotó}}$ és lineal.
\end{itemize}

Si no es distingeixen les dues regions cal restar 0,2 p,

Si no s'indica $W_0$ cal restar 0,1 p,

Si a la primera regió no es representa $E_{C,\text{màx}} = 0$, cal restar 0,1 p,

Si a la segona regió no es representa la dependència lineal, cal restar 0,1 p,

En la representació esquemàtica no és necessari que apareguin les dades del problema
perquè es tracta d'una representació qualitativa. Si l'estudiant intenta fer una
representació tenint en compte els valors numèrics del problema, no se'l penalitzarà. El
que es valora és que sàpiga interpretar i plasmar gràficament el fenomen.
\end{solapartat}
